\documentclass[12pt]{article}

\usepackage[a4paper, margin=1in]{geometry}
\usepackage{amsmath, amssymb}
\usepackage{setspace}
\usepackage{enumitem}

\setstretch{1.15}
\setlength{\parindent}{0pt}
\setlength{\parskip}{6pt}

\begin{document}

\textbf{CSE400 – Fundamentals of Probability in Computing}

\textbf{Lecture 6: Discrete Random Variables, Expectation and Problem Solving}

Instructor: Dhaval Patel, PhD

Date: January 22, 2025

\vspace{1em}

\textbf{1. Random Variables: Motivation and Concept}

\textbf{Definition}

A random variable (RV) is a function that maps outcomes from a sample space to numerical values.

Instead of working directly with outcomes (e.g., coin toss sequences), we work with numbers derived from them.

This allows us to assign probabilities to numerical values.

\textbf{Visualization (as per slides)}

The distribution of a random variable can be visualized using a bar diagram.

The x-axis represents the possible values of the random variable.

The height of the bar at value represents:

Each probability is computed from the probability of the corresponding event in the sample space.

\vspace{0.5em}

\textbf{2. Types of Random Variables}

\textbf{2.1 Discrete Random Variables}

A random variable is discrete if:
\begin{itemize}[leftmargin=2em]
\item It has countable support
\item Probabilities are assigned to single values
\item Each possible value has strictly positive probability
\item It is described using a Probability Mass Function (PMF)
\end{itemize}

\textbf{2.2 Continuous Random Variables (for contrast only)}

\begin{itemize}[leftmargin=2em]
\item Have uncountable support
\item Described using a Probability Density Function (PDF)
\item Probabilities are assigned to intervals
\item Each exact value has zero probability
\end{itemize}

(Note: No continuous RV calculations are performed in this lecture.)

\vspace{0.5em}

\textbf{3. Example 1: Tossing 3 Fair Coins}

\textbf{Experiment Description}

Toss 3 fair coins

Let random variable $Y =$ number of heads observed

$a_i$

$\Pr[X = a_i]$

\textbf{Possible Values of $Y$}

\[
Y \in \{0,1,2,3\}
\]

\textbf{Probability Calculations}

\[
P(Y = 0) = P(t,t,t) = \frac{1}{8}
\]

\[
P(Y = 1) = P(t,t,h) + P(t,h,t) + P(h,t,t) = \frac{3}{8}
\]

\[
P(Y = 2) = P(t,h,h) + P(h,t,h) + P(h,h,t) = \frac{3}{8}
\]

\[
P(Y = 3) = P(h,h,h) = \frac{1}{8}
\]

\textbf{Probability Distribution Summary}

\[
P(Y=0)=\frac{1}{8},\;
P(Y=1)=\frac{3}{8},\;
P(Y=2)=\frac{3}{8},\;
P(Y=3)=\frac{1}{8}
\]

\textbf{Total Probability Check (Assumption: RV must take exactly one value)}

\[
1 = P\left(\bigcup_{i=0}^{3} \{Y=i\}\right)
  = \sum_{i=0}^{3} P(Y=i)
\]

\vspace{0.5em}

\textbf{4. Probability Mass Function (PMF)}

\textbf{Definition}

A random variable that can take at most a countable number of values is called discrete.

Let $X$ be a discrete random variable.

Range (possible values):

\[
R_X = \{x_1, x_2, x_3, \dots\} \quad \text{(finite or countably infinite)}
\]

\textbf{PMF Definition}

The function
\[
P_X(x_k) = P(X = x_k), \quad k = 1,2,3,\dots
\]
is called the Probability Mass Function (PMF) of $X$.

\textbf{Fundamental Property of PMF}

Since $X$ must take one of the possible values:
\[
\sum_{x \in R_X} P_X(x) = 1
\]

\vspace{0.5em}

\textbf{5. PMF Example: Given Functional Form}

\textbf{Problem Statement}

The PMF of a random variable $X$ is:
\[
p(i) = c \frac{\lambda^i}{i!}, \quad i = 0,1,2,\dots
\]
where $\lambda > 0$.

Find:
\begin{enumerate}[leftmargin=2em]
\item $P(X = 0)$
\item $P(X > 2)$
\end{enumerate}

\textbf{Step 1: Find the Constant}

Using the PMF property:
\[
\sum_{i=0}^{\infty} p(i) = 1
\]

Substitute:
\[
c \sum_{i=0}^{\infty} \frac{\lambda^i}{i!} = 1
\]

Using the known series expansion:
\[
\sum_{i=0}^{\infty} \frac{\lambda^i}{i!} = e^{\lambda}
\]

Thus:
\[
c e^{\lambda} = 1 \Rightarrow c = e^{-\lambda}
\]

\textbf{Step 2: Compute}
\[
P(X = 0) = e^{-\lambda}
\]

\textbf{Step 3: Compute}

Using complement:
\[
P(X > 2) = 1 - \left(P(X=0) + P(X=1) + P(X=2)\right)
\]

\vspace{0.5em}

\textbf{6. Bayes’ Theorem (Recap)}

\textbf{Starting Point}
\[
P(A \cap B) = P(B \mid A)P(A)
\]

Using:
\[
P(B_i \mid A) =
\frac{P(A \mid B_i)P(B_i)}
{\sum_j P(A \mid B_j)P(B_j)}
\]

This is called Bayes’ Formula (Proposition 3.1).

\textbf{Definitions}

Prior probability: formed before observing data

Posterior probability: computed after observing event

\vspace{0.5em}

\textbf{7. Bayes’ Theorem – Example 1: Auditorium with 30 Rows}

Auditorium has 30 rows

Row 1 has 11 seats

Row 2 has 12 seats

\ldots

Row 30 has 40 seats

\textbf{Selection Process}

\begin{enumerate}[leftmargin=2em]
\item A row is selected uniformly at random
\item A seat is selected uniformly at random within that row
\end{enumerate}

\textbf{Key Formula Used}

\[
P(S_{15}) = \sum_{k=5}^{30} P(S_{15} \mid R_k)P(R_k)
\]

\[
P(R_k) = \frac{1}{30}
\]

\[
P(S_{15} \mid R_k) =
\frac{1}{\text{number of seats in row } k}
\]

Final numerical value (as shown in slides):
\[
P(S_{15}) = 0.0342
\]

\vspace{0.5em}

\textbf{8. Bayes’ Theorem – Example 2: Communication System (Receiver)}

Binary data is transmitted

Receiver (RX) may make errors:

0 may be detected as 1

1 may be detected as 0

$\{0,1\}$

Given

A set of conditional probabilities describing detection errors

(The lecture slide introduces the setup; detailed numerical evaluation continues beyond this slide.)

\vspace{1em}

\textbf{End of Lecture Scribe}

\vspace{0.5em}

\textbf{Final Notes (for Revision)}

Every PMF must sum to 1

Bayes’ theorem reverses conditioning

All examples strictly follow:

Proper normalization

Complement rule

Law of total probability

\end{document}
